\subtituloA{Limitaciones y seguridad del entorno cliente}

\lista{
  \textbf{Política del mismo origen (SOP):}** Restringe el acceso a recursos basándose en protocolo, dominio y puerto (ej. banco.es no accede a www.otrositio.com). |
  \textbf{CORS:} Error que ocurre al acceder a APIs externas; requiere permiso explícito (cabecera) del servidor externo, no es un fallo del código cliente. |
  \textbf{XSS:} Vulnerabilidad donde texto no limpiado por el usuario puede inyectar y ejecutar código script en la página de otro usuario (Tema 6).
}

\nota{Para el examen}{
Es falso que HTML gestione páginas dinámicas; la gestión dinámica requiere lenguajes de script tanto en el servidor como en el cliente (JavaScript). La limitación clave de JS es la incapacidad de escribir archivos en el servidor.
}